problem:  
Let \(X\) be a compact Kähler manifold with nonnegative holomorphic bisectional curvature, and let its universal cover split holomorphically and isometrically as  
\[
\widetilde{X} \cong \mathbb{C}^q \times X_1 \times \cdots \times X_r,
\]
where each \(X_i\) is an irreducible compact Hermitian symmetric space of compact type and \(\mathbb{C}^q\) is the flat factor (by Mok’s splitting theorem).  

Let \(G\) be a simple algebraic group over a non‑Archimedean local field \(k\), and let \(\mathscr{B}\) be the associated Bruhat–Tits building endowed with its natural CAT(0) structure.  
Suppose \(\rho : \pi_1(X) \rightarrow G(k)\) is a representation such that the induced action of \(\pi_1(X)\) on \(\mathscr{B}\) fixes no point **in \(\mathscr{B}\) or on its boundary at infinity** (i.e., \(\rho\) is nontrivial in the geometric sense, but not assumed to be nonelementary with respect to preserving flats).  

By Gromov–Schoen, there exists an energy-minimizing \(\rho\)-equivariant harmonic map \( f : \widetilde{X} \rightarrow \mathscr{B} \).

**Problem:**  
Assume \(X\) has nonnegative holomorphic bisectional curvature.  

1. Does every such energy-minimizing \(\rho\)-equivariant harmonic map \(f\) have its image contained in a single Euclidean apartment \(A\subset\mathscr{B}\)?  
   Equivalently, does curvature nonnegativity on the source enforce that the induced representation \(\rho(\pi_1(X))\) virtually stabilizes an affine flat in the building?  

2. If the answer to (1) is affirmative, must \(f\) then factor through the flat torus quotient \(\mathrm{Alb}(X) \cong \mathbb{C}^q/\Lambda\)?  
   In that case, does the induced action on the apartment correspond to a virtually abelian subgroup of the affine isometry group \(\mathbb{R}^n \rtimes O(n)\)?  

3. Conversely, if (1) fails, can one construct examples of such \(X\) and \(\rho\) where \(f(\widetilde{X})\) meets several apartments nontrivially—i.e., nonflat harmonic maps—despite nonnegative curvature of \(X\)?  

Put differently, does positive or nonnegative Kähler curvature on the domain obstruct “metric branching” of harmonic images into non‑Archimedean targets, or can non-flat, non-apartment‑contained harmonic maps exist in this curvature setting?

Why is it a "good" problem:  
This version eliminates the logical conflict of the previous formulation by removing the incompatible “nonelementary (no-flat)” assumption, instead focusing on whether curvature forces apartment stabilization—the expected rigidity conclusion itself. The problem now asks whether classical Kähler‑curvature rigidity phenomena persist when the smooth, negatively curved target of Siu–Sampson theory is replaced by a singular, polyhedral CAT(0) building.  

It elegantly bridges differential geometry, non‑Archimedean harmonic analysis, and representation theory. A positive result would parallel Siu-type rigidity in a non‑Archimedean and nonsmooth setting, implying that curvature nonnegativity forces a representation’s image to act affinely on a flat; a counterexample would reveal new geometric mechanisms by which harmonic maps from smooth Kähler manifolds branch across polyhedral flats.  

The statement is conceptually simple—does curvature compel affine flatness?—yet it touches on the analytic foundations of harmonic map theory in singular spaces, the interplay between curvature and group actions, and the geometry of buildings. This makes it a profoundly insightful and cross‑disciplinary research problem.